% このディレクトリで2回実行：pdflatex -interaction=nonstopmode -halt-on-error report.tex
\documentclass[a4paper,10pt,pdflatex,ja=standard,jafont=ipaex-type1,
  magstyle=real,scale=1]{bxjsarticle}
\usepackage{lmodern}
\usepackage{amsmath,amssymb,mathtools}
\usepackage{geometry}
\geometry{reset,left=22mm,right=22mm,top=20mm,bottom=17mm,
  headheight=13pt,headsep=6mm,footskip=10mm}
\usepackage{xcolor}
\usepackage{enumitem}
\usepackage{fancyhdr}
\usepackage{lastpage}

\definecolor{ink}{HTML}{193A5C}
\definecolor{accent}{HTML}{237D82}
\definecolor{soft}{HTML}{EAF3F3}
\definecolor{muted}{HTML}{52616B}
\newcommand{\coursename}{解析学 II}
\newcommand{\handoutnumber}{4}
\pagestyle{fancy}
\fancyhf{}
\fancyhead[L]{\small\color{ink}\coursename{}・レポート課題}
\fancyhead[R]{\small\color{ink}第\handoutnumber 回}
\fancyfoot[C]{\small\color{ink}\thepage/\pageref{LastPage}}
\renewcommand{\headrulewidth}{0.3pt}
\setlength{\parindent}{0pt}
\setlength{\parskip}{2pt}
\setlength{\abovedisplayskip}{5pt}
\setlength{\belowdisplayskip}{5pt}
\setlength{\abovedisplayshortskip}{3pt}
\setlength{\belowdisplayshortskip}{4pt}

\newcommand{\handouttitle}[1]{%
  \begin{center}
    {\Large\mdseries\color{ink}#1}\par\vspace{3mm}
    {\color{accent}\rule{0.88\linewidth}{0.7pt}}
  \end{center}\vspace{-2mm}}
\newcommand{\handout}[2]{%
  \renewcommand{\handoutnumber}{#1}%
  \handouttitle{第#1回\quad #2}}
\newcommand{\topic}[1]{\par\vspace{5pt}%
  {\large\sffamily\mdseries\color{accent}#1}\par\vspace{2pt}}
\newcommand{\keybox}[1]{\par\vspace{3pt}\begingroup
  \setlength{\fboxsep}{4pt}%
  \colorbox{soft}{\parbox{\dimexpr\linewidth-2\fboxsep\relax}{\small #1}}%
  \endgroup\par\vspace{2pt}}
\newenvironment{problems}
  {\begin{enumerate}[leftmargin=2.4em,itemsep=3pt,topsep=3pt,parsep=0pt,
    font=\color{accent}]}
  {\end{enumerate}}
\newcommand{\problemlabel}[1]{{\sffamily\mdseries\color{ink}#1}}
\newcommand{\problemspace}[1][14mm]{\par\vspace{#1}}

\begin{document}
\fontsize{10pt}{13.5pt}\selectfont
\setlength{\abovedisplayskip}{4pt}
\setlength{\belowdisplayskip}{4pt}
\handout{4}{レポート課題（全10題）}


\keybox{範囲：第1回から第4回まで。問題1から9は基礎問題、問題10は発展問題である。\\
答えだけでなく、計算過程や理由も記すこと。不定積分では積分定数を付け、
原始関数を考える区間も明示すること。$\log$ は自然対数を表す。}

\topic{逆関数・逆三角関数と1変数の積分}
\begin{problems}
\item \problemlabel{逆関数。}
区間 $[-2,\infty)$ 上の関数 $f(x)=x^2+4x$ の逆関数 $g$ を求め、
$g$ の定義域と値域を示せ。
\problemspace

\item \problemlabel{逆三角関数の値と微分。}
逆三角関数は授業で定義した主値を用いる。
\par (a) 次の値を求めよ。
$$
 {\rm Sin}^{-1}\left(-\frac12\right),\qquad
 {\rm Cos}^{-1}\left(-\frac{\sqrt3}{2}\right),\qquad
 {\rm Tan}^{-1}\sqrt3.
$$
(b) 次の関数を微分せよ。
$$
 u(x)={\rm Tan}^{-1}(2x)\quad(x\in\mathbb R),\qquad
 v(x)={\rm Sin}^{-1}\left(\frac{x}{2}\right)\quad(-2<x<2).
$$
\problemspace

\item \problemlabel{多項式除法を用いる積分。}
分子を分母で割り、次の不定積分を求めよ。
$$
 \int\frac{x^3}{x^2+4}\,dx.
$$
\problemspace

\item \problemlabel{部分分数分解。}
被積分関数を部分分数に分解し、次の不定積分を求めよ。
$$
 \int\frac{dx}{(x-1)(x+2)}.
$$
\problemspace

\item \problemlabel{平方完成と逆正接。}
分母を平方完成し、次の積分を求めよ。
$$
\quad\int\frac{dx}{x^2-2x+5}.
$$
\problemspace
\end{problems}

\newpage
\handouttitle{第4回\quad レポート課題（つづき）}
\topic{重積分の定義と累次積分}
\begin{problems}
\setcounter{enumi}{5}
\item \problemlabel{リーマン和から重積分へ。}
$K=[0,2]\times[0,1]$、$f(x,y)=x+2y$ とする。
$N,M$ を正の整数とし、$K$ を $x$ 方向に $N$ 等分、$y$ 方向に $M$ 等分する。
各小長方形の右上の点で関数の値を取る。
\par (a) 小長方形の面積と標本点の座標を求め、リーマン和 $S_{N,M}$ を計算せよ。
\par (b) $N,M\to\infty$ の極限から $\displaystyle\iint_K(x+2y)\,dx\,dy$ を求めよ。
\par {\small\color{muted}必要なら $\displaystyle\sum_{i=1}^n i=\frac{n(n+1)}2$
（$n$ は正の整数）を用いてよい。}
\problemspace[12mm]

\item \problemlabel{区分求積による重積分。}
$K=[0,1]\times[0,2]$、$f(x,y)=x^2y$ とする。
$N,M$ を正の整数とし、$K$ を $x$ 方向に $N$ 等分、$y$ 方向に $M$ 等分する。
各小長方形の右上の点を標本点とするリーマン和 $S_{N,M}$ を計算し、
$N,M\to\infty$ の極限として次の重積分を求めよ。
$$
 \iint_K x^2y\,dx\,dy.
$$
{\small\color{muted}必要なら $\displaystyle\sum_{i=1}^n i^2=\frac{n(n+1)(2n+1)}6$
（$n$ は正の整数）を用いてよい。}
\problemspace[12mm]

\item \problemlabel{三角形上の重積分。}
領域 $D=\{(x,y)\in\mathbb R^2:x\ge0,\ y\ge0,\ x+y\le2\}$ を図示せよ。
次の重積分を2通りの積分順序で計算し、値が一致することを確かめよ。
$$
 \iint_D(x+2y)\,dx\,dy.
$$
\problemspace[12mm]

\item \problemlabel{対称性の利用。}
$D=[-1,1]\times[-2,2]$ とする。対称性によって積分が $0$ になる項を理由とともに示し、
次の重積分を求めよ。
$$
 \iint_D\left(xe^{x^2+y^2}+y^2\right)\,dx\,dy.
$$
\problemspace[12mm]

\item \problemlabel{発展：}重積分の歴史を調べよ。誰が発展させたかなどを３００字程度で書くこと。
\problemspace[12mm]
\end{problems}
\end{document}
